\section*{Capítulo \ref{ch:b3:learning-memory} --- Plasticidade neural, aprendizado e memória}

\begin{solution}{exo:b3:learning-memory:1}
A \omterm{def:b3:learning-memory:kinds}{memória declarativa} (fatos, acontecimentos) foi abolida em H.M. --- ele
não conseguia evocar nada de novo --- enquanto a memória não declarativa
sobreviveu: ele aprendeu o desenho em espelho normalmente, sem se lembrar
das sessões. Habilidades: estriado e \omterm{def:b3:nervous-systems:plan}{cerebelo}; medo condicionado:
amígdala; reflexos motores condicionados: \omterm{def:b3:nervous-systems:plan}{cerebelo}; priming: córtex
sensorial.
\end{solution}

\begin{solution}{exo:b3:learning-memory:2}
Quando uma célula pré-sináptica contribui repetidamente para o disparo
de uma pós-sináptica, a conexão entre elas se fortalece. Daí a LTP ser
específica da entrada (só as sinapses ativas, que participaram, mudam),
associativa (uma entrada fraca ativa enquanto a célula dispara sob uma
forte é fortalecida) e cooperativa (entradas em número suficiente
precisam agir juntas para fazer a célula disparar).
\end{solution}

\begin{solution}{exo:b3:learning-memory:3}
Ele só conduz quando o glutamato está ligado (a célula pré-sináptica
disparou) \emph{e} a membrana está despolarizada o bastante para expulsar
o magnésio (a célula pós-sináptica está ativa): as duas condições da
\omterm{def:b3:learning-memory:ltp}{regra de Hebb} numa proteína só. AP5 antes do tétano: nenhuma entrada de
cálcio, nenhuma LTP, embora a transmissão pelos \omterm{prop:b3:learning-memory:nmda}{receptores AMPA} seja
normal. AP5 depois: a LTP já induzida não é afetada --- o receptor é
necessário para a indução, e não para a manutenção.
\end{solution}

\begin{solution}{exo:b3:learning-memory:4}
Curto prazo: modificação covalente de proteínas existentes --- a PKA
fosforila e fecha um canal de potássio, alargando o disparo e aumentando
a liberação; nenhuma proteína nova. Longo prazo: a PKA ativa o CREB,
genes novos são transcritos e novos terminais sinápticos crescem. Um
inibidor da síntese proteica dado durante o treino repetido deixa a
facilitação de curto prazo intacta e abole a memória medida dias depois;
dado depois, nada faz.
\end{solution}

\begin{solution}{exo:b3:learning-memory:5}
$0.7^{n} \le 0.1$: $n \ge 6.5$, portanto $7$ tentativas. A partir de
$0.9\lambda$, $0.9\times 0.7^{n} < 0.1$: de novo $7$ tentativas.
Simétricas porque o mesmo $\alpha$ governa as duas direções. Nos animais
a extinção costuma ser mais lenta e a memória original volta depois de um
descanso (recuperação espontânea): a extinção é aprendizado novo que
suprime o velho, e não seu apagamento.
\end{solution}

\begin{solution}{exo:b3:learning-memory:6}
Autovetores $(1,1)/\sqrt{2}$ com autovalor $1.5$ e $(1,-1)/\sqrt{2}$ com
$0.5$; Hebb gira $\mathbf{w}$ em direção a $(1,1)$, e o neurônio passa a
detectar as duas entradas disparando juntas. Passo de Oja: $y = 1$,
$\Delta\mathbf{w} = 0.1\times 1\times((1,1) - 1\times(1,0)) = (0,0.1)$,
de modo que $\mathbf{w} = (1, 0.1)$.
\end{solution}

\begin{solution}{exo:b3:learning-memory:7}
$500\times 0.05 = 25$ íons livres; $25/6\times 10^{23} = 4.2\times 10^{-23}$ mol em $10^{-16}$ L: \qty{0.42}{\micro mol/L}, quatro vezes o
nível de repouso, mas abaixo do $K_{d}$ da calmodulina --- apenas
ativação parcial; são necessárias algumas aberturas de canal juntas para
alcançar a faixa micromolar que liga a CaMKII.
\end{solution}

\begin{solution}{exo:b3:learning-memory:8}
A é treinada até $V_{A} = \lambda$: a dopamina dispara em salvas nas
primeiras tentativas e silencia à medida que a recompensa passa a ser
prevista. Nas tentativas compostas a predição $V_{A} + V_{B} = \lambda$
já é exata, o erro é zero, a dopamina fica em silêncio e $V_{B}$ nunca
muda: B é bloqueada. Testada sozinha, B não evoca resposta alguma; se uma
recompensa então seguir B sozinha, ela é inesperada e a dopamina dispara
em salva.
\end{solution}

\begin{solution}{exo:b3:learning-memory:9}
(a) Nenhuma LTP na \omterm{def:b3:learning-memory:hippocampus}{CA1} e um déficit de memória espacial, com os demais
aprendizados intactos. (b) Memória de curto prazo normal, memória de
longo prazo ausente em \qty{24}{h}. (c) Nenhum efeito --- a \omterm{def:b3:learning-memory:kinds}{consolidação}
acabou (a menos que a memória seja reativada, quando volta a ficar
lábil). (d) O camundongo se imobiliza no contexto seguro: a memória é
evocada artificialmente pela \omterm{def:b3:learning-memory:hippocampus}{reativação} das células que a codificaram.
\end{solution}

\begin{solution}{exo:b3:learning-memory:10}
$\Delta|\mathbf{w}|^{2} \approx 2\mathbf{w}\cdot\Delta\mathbf{w} = 2\eta
y^{2} \ge 0$: o comprimento nunca diminui e cresce sempre que o neurônio
dispara. Num ponto fixo de Oja, $C\mathbf{w} = (\mathbf{w}^{\mathsf{T}}C
\mathbf{w})\mathbf{w}$, de modo que $\mathbf{w}$ é um autovetor com
$\lambda =
\lambda|\mathbf{w}|^{2}$ e, portanto, $|\mathbf{w}| = 1$. O
crescimento ilimitado saturaria toda sinapse, aboliria a seletividade e
levaria a excitação a disparar sem freio; os normalizadores biológicos
são o escalonamento sináptico (todas as sinapses de um neurônio
reduzidas quando ele está ativo demais), a LTD e a redução das sinapses
durante o sono.
\end{solution}

\begin{solution}{exo:b3:learning-memory:11}
Rato: $0.14\times 3\times 10^{5} \approx 4\times 10^{4}$; humano: $0.14
\times 2\times 10^{6} \approx 3\times 10^{5}$. A estimativa supõe padrões
aleatórios e densos e conectividade total; os padrões reais são esparsos
(o que eleva muito a capacidade) e estruturados, e o \omterm{def:b3:learning-memory:kinds}{hipocampo} é um
depósito temporário que entrega memórias a um córtex de $10^{10}$
neurônios --- o depósito de uma vida é o córtex, e o número hipocampal é
o tamanho de um buffer, e não uma contagem de memórias.
\end{solution}

\begin{solution}{exo:b3:learning-memory:12}
As entradas do olho aberto estão correlacionadas com o disparo da célula
cortical e são fortalecidas; as do olho fechado não estão e, sob
competição por um peso total limitado, enfraquecem --- Hebb mais
normalização. Com os dois olhos fechados nenhuma entrada é favorecida, de
modo que ambas mantêm (fraco) território e o efeito é mais brando. O
período se fecha à medida que os circuitos inibitórios amadurecem, as
redes perineuronais enrijecem em volta das sinapses e as subunidades do
\omterm{prop:b3:learning-memory:nmda}{receptor NMDA} mudam para uma forma menos permissiva.
\end{solution}

\begin{solution}{pb:b3:learning-memory:1}
\textbf{1.} $10\times 10^{3} = 10^{4}$ íons; \qty{2}{\%} livres, $200$;
$200/6\times 10^{23} = 3.3\times 10^{-22}$ mol em $8\times 10^{-17}$ L:
\qty{4.2}{\micro mol/L}.
\textbf{2.} Sim, acima do limiar de \qty{2}{\micro mol/L}. Dois
receptores: $40$ íons livres, \qty{0.83}{\micro mol/L} --- abaixo dele.
\textbf{3.} A \qty{0.83}{\micro mol/L} a calcineurina está ativa e a
CaMKII não: os \omterm{prop:b3:learning-memory:nmda}{receptores AMPA} são retirados e a sinapse enfraquece
(LTD). Um íon, duas enzimas de afinidade diferente: uma elevação grande e
rápida alcança a quinase; uma pequena e lenta, só a fosfatase.
\textbf{4.} Uma segunda entrada dentro de algumas constantes de tempo
soma seu cálcio ao que restou da primeira; depois de \qty{50}{ms} quase
nada resta. A exigência de coincidência do NMDA e o decaimento de
\qty{20}{ms} dão juntos uma janela de uns \qty{20}{ms} para cada lado ---
a janela de tempo de disparo.
\textbf{5.} Os tampões capturam a maior parte do cálcio numa fração de
micrômetro e o pescoço estreito da espinha retarda a difusão para o
dendrito, de modo que a elevação fica na espinha que estava ativa; uma
vizinha a \qty{0.5}{\micro m} nada percebe, e só a sinapse ativa muda.
\textbf{6.} $(70 - 30)/5 = 8$ entradas sincrônicas, ou um potencial de
ação retropropagado do corpo celular: uma sinapse sozinha não consegue
desbloquear seus próprios \omterm{prop:b3:learning-memory:nmda}{receptores NMDA} --- a \omterm{def:b3:structural-biology:allostery}{cooperatividade} está
embutida na física.
\textbf{7.} $(1,1)/\sqrt{2}$ com autovalor $1.6$; $(1,-1)/\sqrt{2}$ com
$0.4$.
\textbf{8.} $C\mathbf{w}_{0} = (1, 0.6)$, $\mathbf{w}_{1} = (1.10,
0.06)$, ângulo com $(1,1)$: $45^{\circ} - 3.1^{\circ} = 41.9^{\circ}$.
$\mathbf{w}_{2} = (1.21, 0.13)$: $38.8^{\circ}$. $\mathbf{w}_{3} = (1.34,
0.22)$: $35.8^{\circ}$.
\textbf{9.} $w_{2}/w_{1} \to 1$ enquanto $|\mathbf{w}|$ cresce sem limite
(por um fator $1.16$ a cada passo).
\textbf{10.} $C\mathbf{w} = (1.28, 1.28)$; $\mathbf{w}^{\mathsf{T}}C
\mathbf{w} = 2.05$; $\Delta\mathbf{w} = 0.1\times(1.28 - 2.05\times 0.8)
= -0.036$ em cada coordenada: $\mathbf{w} = (0.76, 0.76)$, encolhendo em
direção ao vetor unitário $(0.71, 0.71)$.
\textbf{11.} Conjunto: $y = 1.42$; entrada 2 sozinha: $y = 0.71$, a
metade.
\textbf{12.} A entrada fraca dispara enquanto a forte aciona a célula, de
modo que ela fica correlacionada com a saída e a \omterm{def:b3:learning-memory:ltp}{regra de Hebb} a
fortalece --- o neurônio aprende a co-ocorrência.
\textbf{13.} $1 - 0.85^{n}$: $0.56$, $0.80$, $0.96$.
\textbf{14.} $0.85^{n} \le 0.05$: $n \ge 18.5$, portanto $19$ tentativas.
\textbf{15.} Não: a fração da assíntota alcançada após $n$ tentativas não
depende de $\lambda$; a assíntota é que dobra. Uma recompensa maior pode
elevar $\alpha$ (saliência), e um limiar comportamental sobre o valor
\emph{absoluto} é alcançado mais cedo.
\textbf{16.} $V_{B} = 0$: bloqueada. O composto é plenamente previsto, o
erro é zero, e os neurônios dopaminérgicos ficam em silêncio.
\textbf{17.} B sozinha nada prevê, de modo que o erro é $\lambda$ e a
dopamina dispara em salva; $V_{B}$ então sobe como na aquisição,
$\lambda(1 -
0.85^{n})$.
\textbf{18.} O fármaco entrega a salva que significa ``melhor que o
previsto'' depois do que quer que a tenha precedido; a regra incrementa o
valor dessas pistas e ações a cada ocasião, de modo que elas adquirem
valor independentemente de qualquer desfecho real --- a busca se torna
compulsiva.
\textbf{19.} $0.14\times 3\times 10^{5} \approx 4.2\times 10^{4}$.
\textbf{20.} $4.2\times 10^{4}/20 = 2100$ dias, cerca de seis anos; a
\omterm{def:b3:learning-memory:kinds}{consolidação} precisa mover memórias para o córtex e o esquecimento
precisa liberar o depósito.
\textbf{21.} Escrever padrões novos depressa no córtex sobrescreveria as
sinapses que codificam os antigos (interferência catastrófica); uma
\omterm{def:b3:learning-memory:hippocampus}{reativação} lenta e intercalada permite ao córtex integrar o novo com o
velho, enquanto o \omterm{def:b3:learning-memory:kinds}{hipocampo} guarda a memória nova nesse meio-tempo ---
um buffer rápido alimentando um depósito lento.
\textbf{22.} $0.03\times 3\times 10^{5} = 9000$ células por lugar. Os
padrões esparsos se sobrepõem pouco, de modo que muitos mais podem ser
armazenados antes que sua interferência corrompa a recuperação; a
capacidade sobe mais depressa que $N$ à medida que os padrões ficam mais
esparsos, ao custo de cada padrão usar a especificidade de mais células.
\textbf{23.} Compressão de $20$ vezes: células que disparam com
\qty{400}{ms} de intervalo no percurso disparam com \qty{20}{ms} de
intervalo na \omterm{def:b3:learning-memory:hippocampus}{reativação} --- dentro da janela de tempo de disparo. A razão
de ser da compressão é trazer para o alcance hebbiano sequências lentas
demais para serem associadas em tempo real.
\textbf{24.} Com um dia, a memória ainda depende do \omterm{def:b3:learning-memory:kinds}{hipocampo}; com um
mês, ela já foi consolidada no córtex e não precisa mais dele --- exatamente
a perda temporalmente graduada de H.M., com as memórias recentes perdidas
e as remotas poupadas.
\textbf{25.} Cerca de \qty{4.2}{\micro mol/L} de cálcio livre após dez
aberturas; $19$ tentativas para \qty{95}{\%}; uns $4\times 10^{4}$
padrões na \omterm{def:b3:learning-memory:hippocampus}{CA3} do rato.
\end{solution}
